% =====================================================================
%  Chronology of LSSTCam operations in 2025
%  Intended for inclusion in CTN-007.
%
%  Compiled from the Rubin Observatory Slack channels
%    #cam-summit-ops  (top-level summit/camera operations)
%    #cam-operator    (day-to-day operator traffic)
%  covering 2025-01-01 to 2025-12-31 (top-level messages and all
%  thread replies).  All dates and times are UTC.  Ticket identifiers
%  (SUMMIT-, OBS-, FRACAS-, LSSTCAM-, LSSTCCS*-, IT-, DM-, BLOCK-T)
%  are quoted verbatim as they appear in the logs.  Individuals are
%  not named.
%
%  Uses only the standard "description" environment; no extra
%  packages are required.  Temperatures are in degrees Celsius.
% =====================================================================

% Ticket / document link macros (harmless if already defined).
\providecommand{\rubintkt}[1]{\href{https://rubinobs.atlassian.net/browse/#1}{#1}}
\providecommand{\slactkt}[1]{\href{https://jira.slac.stanford.edu/browse/#1}{#1}}
\providecommand{\lsstls}[1]{\href{https://ls.st/#1}{#1}}

\section{Chronology of LSSTCam operations in 2025}
\label{sec:chronology-2025}

The year 2025 took LSSTCam from the Level~3 clean room to routine
on-sky operation on the Simonyi Survey Telescope.  The chronology
below records the milestones of that transition together with the
principal incidents and their resolutions, since much of the
operational context relevant to individual focal-plane channels --
including the R30/Reb1 events discussed in this note -- is only
legible against the wider integration and commissioning history.
It is compiled from the camera operations Slack channels
(\#cam-summit-ops and \#cam-operator); entries are restricted to what
those logs state explicitly.  Throughout, ``Hx'' or ``heat exchanger''
refers to the heat exchanger inside the camera and its vacuum volume, not
to the camera hexapod; where the hexapod itself is meant it is named as
such.

\subsection{January --- final clean-room work and preparation for the lift}
\begin{description}
\item[Jan 5--6] The camera was powered back up in the Level~3 clean
  room after the holiday break, opening a year planned around
  installation, first photon and the start of the survey.  Breakers
  were closed in the sequence main/protection, 24\,V clean, 24\,V and
  48\,V dirty, and it was confirmed that the FES PLC power correctly
  defaults to off.

\item[Jan 7] ComCam was removed from the Top End Assembly (TEA) and
  relocated onto the camera shipping frame, clearing the top end for
  LSSTCam.

\item[Jan 6--15] An intermittent utility-trunk leak-detector fault
  (\texttt{UtLeakDetcFault}) produced three ALARM events and two
  after-hours callouts with no liquid ever found.  A deliberate test
  with Dynalene drops on the leak cable produced the distinct correct
  alert (\texttt{UtCoolantLeak}), showing that the recurring fault was
  a sensor/cable fault; on Jan~15 a load resistor was added to the
  detector's 24\,V-to-5\,V DC-DC converter, after which the fault
  behaved correctly.

\item[Jan 8--14] The TEA was modified to receive the spiral and
  pancake wraps: holes were drilled and tapped through the first
  flange of the camera cable wrap (CCW) with chip-capture tooling and
  the pancake-wrap mounting blocks were installed and torqued.

\item[Jan 7--22] A sustained refrigeration campaign recovered the six
  Level~1 MRCR circuits: filter/dryer assemblies were removed for
  return to SLAC, 2919C refrigerant was recovered from the compressors
  and surge tanks, and all six cryo control boxes were repaired and
  reinstalled.  Open faults included a Cryo~2 pressure transducer,
  Cryo~4 PLC communications and a broken Cryo~5 thermocouple.

\item[Jan 13--16] Dynalene Chiller~02, the unit serving LSSTCam,
  suffered an internal short circuit and was taken out of service;
  Dynalene continued to circulate through the camera via Chiller~01.

\item[Jan 18] Dynalene flow to the camera was lost unexpectedly with
  facility staff off site.  The camera team opened the quadbox
  switches and stopped the FCS, vacuum and shutter processes as a
  precaution; flow was restored by powering the Dynalene cRIO back on.

\item[Jan 10--20] Final filter-exchange-system (FES) and OCS-bridge
  testing was completed, all filters were confirmed stored in the
  carousel, and on Jan~20 the camera and pump cart were powered off at
  the main breakers, removing any dependence on Dynalene.

\item[Jan 20--24] The pancake wrap was installed on the CCW and mated
  to the spiral wrap, every utility was disconnected, the
  utility-trunk skins were installed and the bulkhead limit switches
  were replaced.  Removal of temporary grounding wires exposed an
  unresolved electrostatic grounding problem with the camera doors,
  which ride on non-conductive UHMW tracks; grounding the doors by
  hand was adopted as an interim operational step.

\item[Jan 27--31] Two hardware problems delayed the lift out of the
  clean room.  Torquing the skin fasteners to a newly specified
  34\,in-lb repeatedly stripped screw heads, and several 3/8-16 SHCS
  holding the pancake wrap's central rotating assembly to the adapter
  plate were found sheared, forcing removal and reinstallation of the
  wrap with new hardware at a reduced interim torque.  A separate
  bolt test to 75\,ft-lb showed no failure, leaving the original
  failure unexplained.
\end{description}

\subsection{February --- installation on the Top End Assembly}
\begin{description}
\item[Feb 3--4] The pancake-wrap interface was reworked (pin holes
  reamed, bolts set to half torque, internal sliding blocks removed)
  and the camera was sealed and moved out of the Level~3 clean room.

\item[Feb 5] LSSTCam was lifted from the saddle stand and installed on
  the camera rotator of the Top End Assembly --- the first time the
  camera was integrated with part of the telescope.  Staged hardware
  was found to be the wrong grade and a full set of A4-80 stainless
  bolts was used for the final installation.

\item[Feb 7--11] MTP data fibres at the camera bulkhead were found
  contaminated with Dynalene residue.  The connectors were cleaned and
  Fluke-tested; many failed certification because of scratches, and
  the cables were finally accepted as adequate for the DAQ on a
  best-effort basis (\rubintkt{IT-5920}, \rubintkt{IT-5885}).

\item[Feb 6--13] Summit glycol upgrades and computer-room work left
  the camera unpowered and the control system down for about a week.
  The first power-on of LSSTCam on the TEA on Feb~13 was aborted
  minutes after the breakers closed when a back-flange temperature
  channel ran out of range together with a faint plastic smell; the
  reading was subsequently identified as a bad sensor and the power-on
  was repeated successfully.

\item[Feb 13--14] FES carousel verification failed on the first
  attempt (autochanger truck undervoltage; carousel ``Status 3, error
  register 32'', then a Following Error with no motion) and was
  recovered the following day by reseating the encoder ribbon cable
  between controller and motor.

\item[Feb 17--26] All six cryo circuits were pressurised to 150\,psig
  for pressure-hold tests.  Leaks on circuits 1, 3 and 5 were
  tightened out; two Cryo~2 pressure-transducer channels read
  incorrectly with the fault localised inside the utility-trunk MAQ20
  and worked around using spare cables; a Cryo~5 transducer of the
  wrong pressure range was replaced.  All circuits were declared
  ready for transport by Feb~27.

\item[Feb 18--28] The rotator, camera cable wrap and camera hexapod
  motion-test campaign ran on Level~3, with the camera rotating for
  the first time on Feb~18.  On the same day the CCW local interlock
  at Level~3 was found to be entirely isolated from the rotator
  interlock --- so that the rotator would not stop on a CCW fault ---
  and the two were joined that afternoon.

\item[Feb 21--26] REB power-on and fibre-link verification concluded
  with all 71 camera links declared good; only R20/Reb2 showed link
  errors during the raft-by-raft cycling, and these did not repeat on
  retest.

\item[Feb 21--24] A clogged glycol filter at the Dynalene chiller
  drove utility-trunk Dynalene to $22\,^{\circ}$C over a weekend and
  forced the camera down to minimum power (roughly 400\,W); why the
  alarms did not trigger remained under investigation.

\item[Feb 25--28] A vacuum-subsystem software fault left seven MKS
  gauges --- including the safety-critical cryostat and heat-exchanger
  (Hx) gauges --- offline for over 15\,hours and aborted the low-vacuum
  check, so the \texttt{REFRIG\_NOT\_PERMITTED} alarm was never
  re-raised despite high pressures.  A fix was prepared under
  \slactkt{LSSTCCSSUB-212}.
\end{description}

\subsection{March --- the camera reaches the telescope}
\begin{description}
\item[Mar 1--4] LSSTCam was fully disconnected on Level~3: HFE was
  purged back to the chiller reservoir, a temporary grounding chain
  was built, the CCS subsystems and HCUs were stopped and the camera
  breakers opened.

\item[Mar 3--4] The observatory-wide Cycle~40 update was applied and
  CCS restarted, but the vacuum subsystem crashed on the new build
  with a null-pointer exception in the cryo turbo vent-valve update
  and had to be patched by hand on the summit.

\item[Mar 5--6] The camera was lifted from Level~3 to Level~8 on
  Mar~5 and installed onto the TMA on Mar~6.

\item[Mar 10--13] Camera utilities were reconnected on the TMA
  (Dynalene hoses, MTP fibres, main power, compressed dry air) and
  ComCam was moved back to the white room.  The roughing pump cart was
  swapped after two kinks were found in its hose; a rate-of-rise test
  confirmed the repair.

\item[Mar 11--13] Image acquisition in emulated mode through OCS was
  restored after the DAQ was found to be rejecting triggers because
  the image metadata contained guider sources not present in the
  partition; the emulation configuration was reshaped to science plus
  wavefront sensors only.  Rapid Analysis went live on the summit on
  Mar~13, with camera readouts appearing on RubinTV.

\item[Mar 14--18] The first attempt to power the camera on the TMA
  failed with no network connectivity to any HCU.  The two LC fibres
  of the CCS pair were found swapped relative to the TMA breakout box;
  connectivity was made by changing which media converter carried the
  TMA fibre, and power-on was achieved on Mar~18 (\rubintkt{IT-6016}).

\item[Mar 18--20] The TMA pump cart failed and was replaced by the
  previously removed unit, whose internal gauge was found
  miscalibrated for the summit elevation (corrected in the vacuum
  code).  Both the heat-exchanger and cryostat chambers reached roughing
  vacuum, and repeated USB disconnects of the cart link forced a
  switch to the serial interface.

\item[Mar 18--24] The FES carousel fault carried over from February
  was fixed with a 3D-printed retainer on the encoder cable, and five
  filter exchanges were demonstrated with the camera on the telescope
  on Mar~24.

\item[Mar 24--26] Systematic REB/DAQ link verification found the whole
  focal plane good on the OTM-A side and four bad links on the B side,
  all connected to the same RCE; all four were recovered by swapping
  DPM transceivers in the DAQ, confirming the fault was on the DAQ
  side rather than in the camera or the fibre run.

\item[Mar 27] With the PCS chiller cooled from $+14.5$ to
  $-40\,^{\circ}$C, all REBs were powered on, the GREB firmware was
  updated into guider mode, and roughly 350 full focal-plane images
  with guiding were taken under CCS control.

\item[Mar 28--31] PCS chiller~2 tripped twice when its upstream
  circuit breaker opened; the installed 32\,A breaker was found
  under-rated against the chiller's 50\,A service requirement
  (\rubintkt{OBS-826}).  Restarting the chiller after the breaker work took
  roughly a dozen attempts over two days, repeatedly tripping on low
  HFE flow until the coolant tank was pressurised.  Renewed Dynalene
  instability over Mar~30--31 tripped the REB power supplies on high
  temperature and led the team to conclude that starting the cryo
  system was too risky.

\item[Mar 13, Mar 27] Meeting and shift structure was changed for the
  installation period: a daily Camera Coordination meeting was
  introduced, daily plans moved to threaded summit posts, and a
  weekend on-call person and an overnight alert address were named.
\end{description}

\subsection{April --- first full focal-plane operation and the R30/Reb1 short}
\begin{description}
\item[Apr 1] The cryostat cooldown and refrigeration restart uncovered
  failed water solenoid valves on Cryo~3 and Cryo~6 (swapped with MRCR
  spares) and a Cryo~6 compressor failure; all cryostat and heat-exchanger
  ion pumps were subsequently turned on successfully.

\item[Apr 2] The CCD shorts test passed for nearly the whole focal
  plane.  The exception was R20/Reb2/S21, whose reset-drain voltage
  read an unphysical $+54$\,V instead of the expected value; power
  cycling the REB and resetting the RCE did not change it, and the
  working theory was an open ADC input on the board.  That sensor was
  left off while the rest of the focal plane was powered on.

\item[Apr 3--6] A PCS chiller firmware upgrade was followed by a major
  PCS failure: stage-1 refrigerant leaked out slowly, the control
  valve would not open and flow could not be sustained.  The camera
  was held overnight in a quiescent state with CCDs and REBs off and
  the cryo circuits running, and the chiller was restarted
  successfully on Apr~6 --- described by the team as the fastest
  recovery of its kind.

\item[Apr 7] After a routine power cycle, R30/Reb1 raised a
  \emph{Bias2 OD2 overcurrent} fault on S12, measuring
  $\mathrm{ODI}=0.213$\,A against an $0.082$\,A limit; S12 reached its
  bias voltage in two steps and then tripped off on both attempts.
  The behaviour was assessed in the moment as consistent with a real
  short.  All CCDs were powered off camera-wide and further CCD
  power-on was suspended pending expert review, in part because the
  fault appeared after that day's reorientation of the camera to
  zenith.  R30/Reb1 was left powered off.

\item[Apr 7--8] Repeated CCS jGroups multicast storms disrupted
  trending and control, apparently triggered by starting CCS processes
  on the spare shutter subsystem whose network port lacked IGMP
  snooping; the summit Cycle~41 software and OS upgrade fell on the
  same day.

\item[Apr 8--11] Intermittent PCS chiller~2 trips with no
  corresponding chiller anomaly were traced to an unseated pin in a
  refrigeration-permit connector on the telescope-side cabling,
  aggravated by TMA shaking tests; the pins and shell were replaced
  (\rubintkt{SUMMIT-10642}, \rubintkt{SUMMIT-10652}).

\item[Apr 11] The L1 lens cover and lens cap were removed --- a step
  towards first photon --- and a shutter and CCD soak test was
  authorised with the dome lights administratively controlled.

\item[Apr 13--25] The filter carousel developed an intermittent
  communication fault on socket~3 that depended on rotation speed and
  telescope orientation.  It was diagnosed as a sticking spring-loaded
  carbon collector contact rather than a wiring problem, and the
  collector was replaced on Apr~24--25 (\rubintkt{SUMMIT-10680}), with the
  ageing of the other eight collectors noted as a residual concern.

\item[Apr 14--16] The first full focal-plane CCD power-on campaign was
  run raft by raft, followed by the first HV bias ramp across the
  REBs, which took about 1.5\,hours to reach full voltage.  R30/Reb1
  again tripped on an OD2 overcurrent (0.087\,A) and was left off with
  R20/Reb2.  Over the same nights the camera took part in TMA soak
  testing (BLOCK-T447, BLOCK-T449), giving tracking with shutter
  readouts for the first time, and in the first bright-star tests,
  whose largest observed CCD current jump was benign.

\item[Apr 16--19] Guider idle-clearing was found to inject noise into
  neighbouring sensors during bias readout, making early bias images
  unusable; manual guider sleep/wake was adopted as an interim
  workaround while a proper non-guider clearing scheme was developed
  (\slactkt{LSSTCCSRAFTS-778}).

\item[Apr 17--29] A series of HV bias trips was traced to stray light
  reaching the focal plane: head lamps and light past the filter-loader
  door during FES work, sunlight through an open P-flow door with the
  mirror covers open, and a laser tracker left running.  New
  procedures were proposed tying HV state to P-flow and mirror-cover
  status.

\item[Apr 24] R20/Reb2 was returned to service by configuring the
  erroneous reset-drain reading to be ignored in software rather than
  repairing the board, and HV was restored to the full camera.  The
  same day a REB-firmware register-size bug was fixed, allowing three
  previously off corner CCDs to be powered and validated.

\item[Apr 24--25] A renewed attempt to recover R30/Reb1 was abandoned.
  A shorts test showed an unexplained digital-current spike of order
  1\,A on the REB, seen in neither repeated retests nor 24\,hours of
  monitoring of any other REB, and it was decided not to attempt
  powering it on; R30/Reb1 ended the month as the only REB switched
  off camera-wide.

\item[Apr 26--28] A new sequencer file was tested on sky
  (BLOCK-T469, BLOCK-T471) to address e2v noise and covariance
  anomalies; e2v channels improved clearly while ITL channels changed
  little, and the default sequencer was restored for scheduled
  active-optics measurements.
\end{description}

\subsection{May --- cold weather and the first loss of cryo capacity}
\begin{description}
\item[May 1--7] Memory faults on the data-collection nodes forced
  image handling off \texttt{lsstcam-dc10} onto \texttt{lsstcam-dc02},
  with a DIMM replacement on another node and a full reboot of the
  LSSTCam machines on May~7 with the camera taken offline first.

\item[May 4] A filter-repeatability test using the collimated beam
  projector (BLOCK-T486) measured spot shifts of roughly 5\,pixels
  (about 52\,$\mu$m) before and after a filter swap, within the
  100\,$\mu$m repeatability specification.

\item[May 6--20] Four in-camera filter exchanges were carried out at
  horizon pointing, taking the loaded complement through
  $[g,r,i,z,y]$ to $[g,r,i,z,u]$.  Extraction of a single filter took
  slightly under an hour; swaps required the mirror cover closed
  against light leaks and the HV switches open.

\item[May 7] The cryostat pressure reached the $10^{-8}$\,Torr regime
  for the first time in the reported monitoring.

\item[May 8--11] The camera lost essentially all cryo cooling capacity
  overnight (\rubintkt{OBS-931}, \rubintkt{FRACAS-280}).  The cryo plate began deviating
  from $-127\,^{\circ}$C in the early hours, all five running circuits
  appeared to lose capacity, and CCDs and REBs were shut down around
  08:10--08:20 UTC as both cryostat and heat-exchanger vacuum climbed; the
  heat-exchanger gauge reached about $1\times10^{-5}$\,Torr against an
  ion-pump shutdown threshold of $1.1\times10^{-5}$\,Torr.  Oil was
  found dripping from a liquid-line filter dryer, three dryers were
  swapped, and recovery took three days, with all REBs back on by
  May~10 and HV restored the same evening.  Low ambient temperature,
  with a transition near $7\,^{\circ}$C, was identified as the likely
  driver: cryo trim heat went from about 60\,W before the event to
  about 180\,W after recovery.

\item[May 9--11] The focal plane was fully re-powered after the
  cooling event, with shorts tests on all sensors and REB power-on
  gated on power-supply board temperature.  R30/Reb1 was deliberately
  skipped, its HV switches left open by direction; R21/Reb0, which had
  repeatedly glitched near 50\,V, was set to run permanently at
  46\,V.

\item[May 12--27] A cold-weather operating regime was established.
  Quantitative ``leading indicator'' thresholds (Cryo~1 heat-exchanger return
  temperature below the dome dew point, trim-heat drift, dome ambient
  below $8\,^{\circ}$C, rapid warming of circuit exit temperatures)
  were defined and put on Chronograf dashboards for the observing
  specialists, later reworked to use derivatives; the dome-warming
  procedure was rewritten to use HVAC only.  Two further near-crashes
  (May~16--17 and May~25--27) were managed by closing the dome,
  running closed-dome calibrations and walking the telescope down to
  low elevation to keep the Cryo~1 heat-exchanger return above the dew point.

\item[May 19--23] A refrigerant-additive and insulation campaign
  recovered a large part of the lost capacity: R124 and R245fa were
  added to four circuits and 2919C to all five running circuits, after
  which cryo capacity was assessed as roughly 90--100\,W higher, and
  the Level~8 cryo liquid supply lines were progressively insulated.

\item[May 21] A power glitch stuck the glycol remote control valve
  closed, dumping the glycol charge and taking down the PCS and three
  of four running cryo circuits (\rubintkt{FRACAS-282}).  With the CCS shell
  effectively unusable, REBs were powered off with a direct command;
  the cold plate rose to $-16\,^{\circ}$C before flow was restored via
  the manual valve.

\item[May 15--16] Two software faults cost observing time: a 0.1\,s
  exposure request left the focal plane stuck integrating because the
  configured minimum exposure time was below the shutter's real 1\,s
  minimum (\rubintkt{OBS-947}), and a broadcast packet storm from three camera
  hosts faulted the camera (\rubintkt{OBS-950}).

\item[May 28--30] A new CCS refrigeration release added interlocks and
  alarms drawn from the month's failures (blocking coolant flow below
  a configurable glycol flow, warning on high discharge and low return
  pressure, fixing reversed Cryo~3 water connections), and approval was
  given to replace the failed Cryo~6 module with the Cryo~5 unit from
  the MRCR.
\end{description}

\subsection{June --- First Look, and completion of the science array}
\begin{description}
\item[Jun 1--4] Throttling the Cryo~1 liquid-line isolation valve
  rebalanced liquid and vapour flow and, for the first time, allowed
  Cryo~1 to run continuously overnight at a workable heat-exchanger return
  temperature.

\item[Jun 2--3] The shutter jammed part-open in the middle of a flat
  calibration and could not be recovered by the standard procedure;
  the blade appeared to overrun its high limit until both PLCs and the
  HCU were power-cycled and the shutter recalibrated (\rubintkt{OBS-1013}).

\item[Jun 3--4] Dynalene flow was lost overnight, forcing an
  emergency power-down of CCDs, REBs and most electronics and a switch
  to the second Dynalene chiller (\rubintkt{OBS-1014}).  During recovery HV was
  enabled with the REB back-bias switches still open, producing strong
  tree rings and large gain shifts in flats; this was corrected the
  next morning by zeroing HV and closing the switches raft by raft.

\item[Jun 2--7] The Cryo~6 compressor module was replaced with the
  MRCR unit, and the circuit was evacuated, charged, started and
  gas-chromatographed, restoring six of six cryo circuits.

\item[Jun 4--11] Recurring ``timed out waiting for shutter exposure to
  complete'' failures were traced to PTP clocking rather than to the
  shutter mechanism.  Enabling a boundary clock on the shutter's
  network switch reproduced the fault on the clean-room system;
  transparent-clock mode and then the leaf switch as boundary clock
  eliminated the offset from master.  The investigation is documented
  in \href{https://ctn-003.lsst.io}{CTN-003}.

\item[Jun 6--7, Jun 16--18, Jun 27--30] Three severe cold-driven cryo
  degradation events dominated the second half of the month.  With
  dome air down to $3.6\,^{\circ}$C, then near $5\,^{\circ}$C, then
  near freezing, cryo-plate trim heat was driven to zero, and each
  event required powering off the CCDs, the HV and --- in the second
  and third events --- nearly all REBs, with recovery taking from
  seven hours to three days.  Moving the telescope from high elevation
  to horizon was found to produce a step recovery in cooling capacity.

\item[Jun 7] A multicast storm on the camera network, attributed to
  enabling IGMP snooping on one switch, forced most CCS services to be
  stopped and restarted and left the camera in fault.

\item[Jun 10--11] A screw dropped into the camera body through the FES
  access during Auto Changer PLC work (\rubintkt{SUMMIT-10747}) triggered a
  stop-work and cancelled that night's observing; the camera was
  wrapped, scaffolding erected and side panels removed, and the screw
  was recovered the following morning.

\item[Jun 12--17] Cold-sensitivity mitigation was attempted with
  heater blankets and insulation on the refrigeration lines --- a
  thermocouple measurement initially showed the existing insulation
  had ``zero effect'' because the lines were bare up- and downstream
  --- and refrigerant was removed from Cryo~1 to raise its
  heat-exchanger return temperature clear of the dew point.  Of four dome air
  handlers only one was healthy, so dome heaters were ruled out as a
  mitigation.

\item[Jun 14--15] The compressed dry air supply to the Top End
  Assembly was found to have been silently off for an unknown period,
  and was restored; the TEA air valve was shown to have little
  authority over humidity in the camera region, and the camera team
  stated that dry air to the camera was necessary to survive the
  winter.

\item[Jun 23] Rubin Observatory's public First Look event took place.
  Camera shifts were flexed the previous night around it, and a
  telescope slew with the mirror covers open and dome floodlights on
  during the associated demonstration tripped HV on one REB.

\item[Jun 23--30] CCS 41.0.4 was activated in a staged full-summit
  restart, delivering idle-clear support and the software needed to
  turn on the remaining R30 CCDs.  Idle clear was tested on sky but
  initially kept disabled because early measurements suggested it
  raised the read noise; after two locking bugs were fixed --- skipped
  invocations piling up and a sequencer-reload race --- it was
  declared available for operational use on Jun~30 and nightly
  operations moved to the 3\,s readout sequencer with idle clear
  enabled.

\item[Jun 25--26] The last two focal-plane sensors, on R30/Reb1, were
  powered on and biased, completing the science array.  The turn-on
  required an R30/Reb1 firmware update, an NTP fix, and configuration
  changes that hit a hard-coded firmware parameter range, forcing an
  emergency change to the bias-control limits and a same-day rebuild
  of the rafts subsystem.  During the subsequent telemetry review the
  CCD current on S12/Seg17 was noted as anomalously low
  ($0.017$), and the channel was again described as the only
  non-nominal one when R30/Reb1 was checked on Jun~29.

\item[Jun 28--30] During re-power-up after the third cold event,
  R30/Reb2 showed its clock-high current ramping continuously for
  about an hour, prompting a near stop-work.  The cause was found to be
  idle clear interacting with the REB because the idle-clear pause had
  timed out during the shorts test; a power cycle with idle clear
  correctly paused resolved it, and disabling idle clear during REB,
  CCD and HV turn-on became standing practice.
\end{description}

\subsection{July --- operational maturity of idle clear and the FES}
\begin{description}
\item[Jul 1--4] Idle clear was validated through a series of
  diagnostic bias and dark acquisitions and adopted as the default
  configuration for all on-sky observations, replacing continuous
  all-day bias acquisition.

\item[Jul 1, Jul 15] Filter swaps continued: $u$ out and the pinhole
  filter in on Jul~1 (followed by an approved pinhole illumination
  test with HV off), and $y$ out and $u$ in on Jul~15
  (\rubintkt{SUMMIT-10923}).

\item[Jul 4] A power outage briefly interrupted both glycol and
  Dynalene flow, and glycol chiller~3 showed a rising-temperature
  anomaly overnight (\rubintkt{OBS-1074}), prompting discussion of resilience to
  future outages.

\item[Jul 5--11] A formal shutter and camera-to-OCS-control training
  programme for observing specialists and commissioning scientists was
  launched, together with an overhaul of the CCS command-level policy
  that granted a trained-observer group level-1 access to the filter
  changer and shutter subsystems.

\item[Jul 6--13] Recurring carousel clamp faults, which had been
  costing significant observing time, were addressed with an automated
  unclamping-recovery patch; by Jul~13 recovery was succeeding
  essentially every time, sharply reducing the need for expert
  intervention.  On Jul~20--21 the FES passed 1000 filter changes
  since April with the error rate below 1 in 400.

\item[Jul 11--13] A dangling guider-initialisation command left the
  guider stuck in a paused state for as long as 22\,hours undetected,
  halting guider file production; the root cause was traced to
  double-commanding between CCS and OCS.

\item[Jul 13--28] Repeated \texttt{ocs-bridge} failures to release the
  mount and rotator motion locks after filter changes cost observing
  time on several nights (\rubintkt{OBS-1098}, \rubintkt{OBS-1079}, \rubintkt{OBS-1119}), tentatively
  linked to a SAL/Kafka interface version mismatch and worked around
  by manual restarts.

\item[Jul 15--16] Candidate REB FPGA firmware was loaded into the
  inactive EEPROM slot on all 71 xREBs in subsets --- wavefront and
  guider REBs, R30/Reb1, and the remainder --- in preparation for a
  later activation during a planned power cycle.

\item[Jul 23--29] A serious clamp fault developed at carousel
  socket~3, with repeated position-sensor errors and failed locking
  recovery suspected to be an autochanger alignment problem.  The
  socket was taken out of service and the $g$-band filter was removed
  from the camera on Jul~29 for bench inspection.

\item[Jul 16--31] Cold-weather mitigation continued with heater tape
  and insulation on the cryo liquid lines, refrigerant top-ups on
  Cryo~1, 2 and 6, and the adoption of quantitative dome open/close
  temperature criteria derived from analysis of the four earlier cryo
  crashes.
\end{description}

\subsection{August --- winter storms, power failures and firmware}
\begin{description}
\item[Aug 1--3] A winter storm closed the mountain road with outside
  temperatures down to $-5\,^{\circ}$C; most of the crew descended
  before the closure and monitored remotely while Cryo~1 and Cryo~2
  lost some capacity.

\item[Aug 5--12] Filter operations were disrupted twice: the $u \to y$
  swap was rescheduled after a clamp timeout during a new
  interface-cleaning step and completed on Aug~7 (\rubintkt{SUMMIT-11029}), and a
  defective batch of filter pins was found to trip FCS limit errors
  and was replaced before the $g$ filter was reinstalled
  (\rubintkt{SUMMIT-11092}, \lsstls{LCA-13232}).

\item[Aug 4--12] A repetitive clicking picked up by the shutter
  microphone, described as sounding like a Geiger counter, was analysed
  for the first time.  It ran from 23:06:13 to 23:07:18\,UTC on Aug~4 at
  2--6\,kHz while the telescope was slewing; M2-hexapod telemetry showed
  the hexapod holding position with its actuator currents oscillating,
  which made it the first suspect.  On Aug~12 the flapping straps of a
  temporary safety barrier on the mesh platform, blown by the HVAC fans,
  were proposed instead and taped down.

\item[Aug 11--16] Power glitches during summit UPS and generator
  inspection tripped the PCS chiller twice on a drive fault
  (\rubintkt{FRACAS-292}) and forced an emergency camera shutdown.  The outage was
  used to boot the xREB EEPROMs into the candidate firmware slot, which
  preserves HV bias across REB resets.  Recovery was then complicated
  by days of jGroups retransmit storms that broke CCS--DAQ
  communication whenever an image was taken with the guider running;
  by Aug~15 the trigger was traced to guiding at an untested 5\,s
  exposure cadence combined with high DAQ node traffic, and reverting
  to the normal cadence stopped the storms.

\item[Aug 15] The cryo system was deliberately driven towards its
  crash threshold, with the dome air handler stopped and the glycol
  chiller lowered, to validate a lower glycol setpoint.  The system
  held --- camera air temperature reached about $5.7\,^{\circ}$C with
  trim heat above 90\,W --- and the test was curtailed only because of
  a condensation risk.

\item[Aug 18] A cold night pushed ambient temperature more than two
  degrees below the phase-separator setpoint, causing refrigerant
  condensation in the vapour lines of all six cryo circuits and
  warming the cryostat plate; the glycol chiller was lowered further
  and a portable 60\,kW heater was deployed as mitigation.

\item[Aug 21--23] CCDs had been running measurably warmer since the
  Aug~11 firmware change.  Power-cycling the whole focal plane and
  rolling the firmware back confirmed a genuine power-draw regression
  of roughly 30\,W, tentatively attributed to a new sequencer-override
  register; the apparent loss of cryo capacity during the storm had in
  part been this self-heating.  R30/Reb1 was left running its own
  firmware version, required for its bias-threshold adjustment.

\item[Aug 23--27] A new automated HV-bias ramp-up and regulation
  subsystem ran through a full night on sky for the first time and was
  then enabled camera-wide, with the procedure trained to a new
  operator.  An anomalously high HV current on R03/Reb1 was managed by
  temporarily lowering its setpoint (\rubintkt{LSSTCAM-124}, \rubintkt{LSSTCAM-125}) and
  resolved by Aug~26.

\item[Aug 23--30] Both Level-1 air compressors degraded in the cold
  --- the dryers froze at $-8\,^{\circ}$C, blocking airflow --- and
  were fitted with insulating jackets.  Loss of the dry-air purge was
  a direct camera concern, both for the vacuum gate valves and for
  frosting on the L3 lens.

\item[Aug 26--27] Large-scale summit OS and Kubernetes patching
  (\rubintkt{SUMMIT-11198}) forced a reboot cascade across the camera servers.  A
  scripting error included the shutter HCU, which had been meant to be
  excluded, and an authentication race condition then cost access to
  it for several hours; the incident was judged serious because the
  same fault on a hardware-safety host would have been far worse.

\item[Aug 30--31] The summit generator tripped, cutting glycol flow
  and leaving the camera on UPS alone.  REBs and their power supplies
  were deliberately powered off while cryostat pressure rose towards
  the ion-pump shutdown limit; recovery of the refrigeration, shutter,
  FES and DAQ took the rest of the day, and the full camera power-on
  was completed on Aug~31 after overnight snow.
\end{description}

\subsection{September --- recovery, then preparation for the shutdown}
\begin{description}
\item[Sep 1--5] The camera recovered from the generator failure, whose
  root cause was ice and snow blocking the generator air inlet
  (\rubintkt{FRACAS-299}); the summit server-room UPS had separately exhausted its
  roughly 30-minute battery and dropped camera network communications
  (\rubintkt{FRACAS-300}).  Recurring generator voltage glitches caused repeated
  minute-scale glycol interruptions --- close to the observed 90 to
  120\,s PCS trip threshold --- and observing on generator power was
  restricted in pointing.  Commercial power was confirmed restored on
  Sep~5.

\item[Sep 3--9] Cryo~1 again showed cold-ambient degradation on
  several nights as outside temperatures approached the
  phase-separator setpoint.  It was decided that the circuit need not
  be switched off, and Cryo~1 was removed from the dome-closure
  decision criteria.

\item[Sep 6] An HV trip occurred with the moon only about 20\,degrees
  away, and after recovery the guider was left in an inconsistent
  state between CCS and the DAQ; the case prompted discussion of a
  proper fault path for HV bias trips at the control-system level.

\item[Sep 7] The same sound returned over the shutter microphone during
  20-degree camera rotations every three minutes, which disposed of the
  safety-barrier explanation of August: this time the M2 hexapod was
  stationary while the camera hexapod moved with every rotation, though
  M2-hexapod motor currents still spiked at each one.  The recirculating
  ball bearings of the hexapod roller-screw actuators, or of the rotator,
  were proposed as the source and a ticket was opened
  (\rubintkt{OBS-1158}).

\item[Sep 14--18] A cluster of DAQ and control-system faults occurred:
  a data-transport module on one crate needed rebooting for the second
  time in a month (\rubintkt{LSSTCAM-130}), leaving the DAQ data store
  inconsistent and image acquisition failing; a storage-node switch
  failed; and a duplicated guider-initialisation left the system in an
  error state.

\item[Sep 17--25] The refurbished vacuum pump cart was commissioned on
  the camera elevation platform, and both the cryostat and heat-exchanger
  turbo pumps were run for their semi-annual maintenance cycle ahead
  of the October cryo work.  Entanglement of the cart hoses during
  platform motion was flagged as a safety hazard.

\item[Sep 20--21] Cryo-line heater-tape circuits repeatedly blew
  fuses after termination work --- five to six failures across several
  attempts, including with slow-blow replacements --- and the cause
  was not resolved within the month.

\item[Sep 21--22] Dome lights were switched on briefly by accident,
  tripping roughly 20 REB HV bias channels, all on ITL rafts;
  recovery followed the documented HV turn-on procedure without
  incident.

\item[Sep 23--25] The $r$ and $g$ filters were extracted ahead of the
  October carousel maintenance (\rubintkt{SUMMIT-11343}), leaving $i$, $z$ and
  the pinhole filter in the camera, and setting a record of
  2\,h\,10\,min for a two-filter operation.

\item[Sep 9--30] The camera state for the coming shutdown was defined
  and executed --- REBs on for thermal monitoring, CCDs off, HV off
  --- while the CCS database underwent a raw-data partitioning
  migration and CCS Cycle~41 was rolled out to all summit nodes.
\end{description}

\subsection{October --- the maintenance shutdown and the loss of Cryo 3}
\begin{description}
\item[Oct 1] Starting a new turbo-test-station subsystem broadcast a
  ``refrigeration not permitted'' message that the refrigeration
  subsystem accepted as real, shutting down all six cryo cabinets.
  Three circuits then refused to restart on low discharge temperature
  and needed a heat gun on the lines before they would run.

\item[Oct 1--20] A roughly three-week planned camera downtime ran in
  parallel with facility work: the FES team serviced and reinstalled
  the carousel clamps on sockets 1--4 and fixed an autochanger
  overcurrent problem, cryo filter dryers were replaced on three
  circuits, CCS was upgraded to cycle~42, the MCM, DAQ-management and
  diagnostic-cluster machines were replaced, and the HCU operating
  systems were patched.

\item[Oct 3, Oct 9] The camera cable wrap interlock faulted and
  blocked rotator motion during FES maintenance rotations, recovered by
  bypassing the interlock at the mount engineering interface and
  re-aligning the wrap to the rotator; the same problem recurred on
  Oct~9 and cost the FES team over an hour.

\item[Oct 3] A major fibre-optic failure on the AURA base to Pachón
  link, ongoing since Sep~30, left base-to-summit communications
  fluctuating between the primary fibre and a microwave backup, with
  no estimated repair time; camera systems were not reported as
  directly affected.

\item[Oct 14--15] The $+X$ deployable platform jammed about 42\,cm
  short of full extension after a hose-drape test, the extension
  encoder having mis-reported position because of motor slip under
  excess load.  Both platforms were taken out of service and returned
  the next day under new hard load limits with no personnel aboard
  during motion.

\item[Oct 15--20] Cryo~3 developed a capillary or expansion-valve clog
  during recharge that could not be cleared by power cycling,
  refrigerant addition or running the telescope horizontal.  The
  circuit was given up on Oct~20 and the camera continued on five of
  six circuits, with the cryo trim setpoint raised to recover about
  10\,W of margin; procedure revisions on isolation-valve handling
  were planned in response (\lsstls{LCA-14625}, \lsstls{LCA-17726}).

\item[Oct 13--16] Candidate REB, WREB and GREB firmware was loaded to
  the spare EEPROM slot and booted across the focal plane, including
  the special R30/Reb1 variant.  The rollout repeatedly exposed a REB
  serial-number read race that left REBs in an invalid validation
  state, and a sequencer opcode-error bug when REBs were power-cycled
  while idle clear was active, mitigated by delaying the start of idle
  clear until the subsystem reached its operational state.

\item[Oct 20--21] The camera returned to on-sky operation.  The first
  night surfaced a stuck shutter requiring reinitialisation
  (\rubintkt{OBS-1013}), incorrect guider bias voltages in the new guider
  firmware, which was rolled back, a recurrence of a long-standing
  wavefront-REB power-on issue (\rubintkt{LSSTCAM-37}), and corrupted FITS
  headers introduced by a recent software change --- the last fixed and
  verified end to end the same day.

\item[Oct 22--23] The final science-filter complement was loaded via
  the FES after re-qualification at horizon and zenith, and a new rule
  was deployed requiring the FES to be in a loaded state before it can
  be handed to observatory control.  With cryo capacity reduced by the
  loss of Cryo~3, the CCD and raft-support-assembly setpoints were
  raised by about one degree for thermal margin.

\item[Oct 28--30] Two Dynalene cooling-loss events forced protective
  shutdowns.  In the first, REB board temperatures approached the
  HV-trip limit and REBs were powered off; on power-up they reported a
  zero serial number and stayed invalid, needing repeated RCE resets
  and clearing of a DAQ opcode-error register.  The second, caused by a
  crash of the Dynalene control computer, required a full camera
  shutdown and a two-to-three hour recovery.  A review of the earlier
  Aug~11 event had established that the camera tolerates only about
  13\,minutes without Dynalene flow before board temperatures approach
  the trip limit.
\end{description}

\subsection{November --- control-system instabilities and the Cryo 6 trip}
\begin{description}
\item[Nov 1--3] A subsystem network fault pegged the REB-power process
  at full CPU and filled the management host's filesystem, and on
  Nov~3 a major jGroups meltdown during a dark sequence forced nearly
  every subsystem to be killed and restarted.  In the aftermath the
  shutter's PTP synchronisation fell badly out of lock and required
  repeated PLC power cycles and a full recalibration in a closed dome
  (\rubintkt{LSSTCAM-144}; recovery procedure in \href{https://ctn-003.lsst.io}{CTN-003}).

\item[Nov 4] A planned swap of the Dynalene control computer required
  a roughly three-hour full camera power cycle, with REBs powered off
  in groups to protect the power-supply boards while flow was cut.
  Several REBs returned in an invalid state and a corner REB tripped
  because its heater had not been switched off first.

\item[Nov 5] A power glitch with no UPS backing tripped the PCS
  chiller and the Cryo~6 compressor and briefly dropped Dynalene flow;
  systems recovered within about ten minutes, and a second smaller
  event later the same day again disturbed the chillers.

\item[Nov 6] A dynamic guider region-of-interest test crashed the
  focal-plane subsystem: an out-of-range ROI produced a parameter
  exception and a core dump rather than a rejected command, and
  recurred on a second attempt (\slactkt{LSSTCCSRAFTS-821}).

\item[Nov 10--21] Weekly planning surfaced thermal problems in the Top
  End Assembly --- infrared imaging showed hot spots near the shutter
  electronics box and cavity humidity above ambient, attributed to a
  dry-air leak near the heat-exchanger return --- and the FES carousel
  developed an intermittent following-error and overcurrent fault on
  the first move of the night, addressed incrementally with a warm-up
  brake sequence and automated recovery.

\item[Nov 17--18] A scheduled summit OS and Kubernetes upgrade took
  the control system offline for a morning, after which the camera
  team rebooted essentially every LSSTCam server and verified the DAQ
  link speeds before returning the camera to observatory control.

\item[Nov 18--19] Hours after the Top End Assembly air nozzle was
  replaced, Cryo~6 was lost: the circuit's PLC was reported dead, the
  discharge temperature spiked to $122\,^{\circ}$C and the oil level
  read negative.  The PLC recovered by itself after about half an hour.
  The root cause was traced to a software update elsewhere on the
  summit mis-setting the chiller setpoint in response to bad
  environmental temperature readings, and the control logic that closes
  the coolant valve on a dead PLC was flagged for revision
  (\rubintkt{FRACAS-322}, \rubintkt{OBS-1394}, \rubintkt{OBS-1395}).

\item[Nov 20] A generator-to-commercial power switchover caused a
  total loss of Dynalene, and the automatic REB shutdown response
  failed because the board-temperature limit had been left set too
  high.  The team shut the camera down manually and carried out a
  multi-hour recovery including RCE reboots, invalid REB states and a
  shutter PTP resynchronisation (\rubintkt{FRACAS-321}).

\item[Nov 21--24] The focal-plane subsystem repeatedly hung in an
  active command state after guider initialisation on several nights,
  once timing out after 45\,minutes and each time requiring a subsystem
  restart (\rubintkt{LSSTCAM-148}); a parallel memory-exhaustion event in the
  message broker triggered an overnight jGroups storm (\rubintkt{LSSTCAM-149},
  \rubintkt{LSSTCAM-150}).  On Nov~24 the DAQ link on one data-collection node was
  found to be dropping packets, with register-read retries up by one to
  two orders of magnitude since Nov~20, and the link was bounced and
  the switch power-cycled (\rubintkt{LSSTCAM-151}).

\item[Nov 24--25] A newly trained loader operator completed a filter
  swap using the simplified loader procedure without camera-team
  intervention for the first time, following weeks of FES software
  fixes.

\item[Nov 24] The ``natural frequency'' project, retuning the REB
  sequencer timing to 6.4\,ns, neared completion and was revalidated on
  the test stand, with dedicated dark-dome engineering time requested
  ahead of a deployment decision.
\end{description}

\subsection{December --- refrigeration margin and firmware campaigns}
\begin{description}
\item[Dec 3--11] With Cryo~3 out of service and failing to restart, a
  capillary clog developed on Cryo~4 and drove cryo-plate trim heat
  down to about 1\,W whenever the telescope moved to high elevation.
  The raft-support and cryo-plate setpoints were raised in steps as
  mitigation; the Cryo~4 clog cleared spontaneously on the evening of
  Dec~10, trim heat recovered to tens of watts, and the setpoints were
  restored.  An average cryostat temperature of $-120\,^{\circ}$C was
  documented as the hard safety boundary against sublimation.

\item[Dec 8--9, Dec 13] Two FES faults required expert intervention:
  repeated PLC faults on an online clamp overnight on Dec~7--8, and
  autochanger truck controllers going to fault with undervoltage
  errors after the carousel was found out of standby, reset remotely
  (\rubintkt{OBS-1454}).

\item[Dec 9, Dec 14--15] The 6.4\,ns natural-frequency REB firmware was
  tested on the summit in two phases, booted from the spare EEPROM
  slot.  Both attempts left several REBs hung or in an invalid state
  after RCE resets, and the planned bias and flat sequences
  (BLOCK-T659, BLOCK-T661) were only partly completed before the
  production firmware was restored.

\item[Dec 15--16] Water droplets and streaks were found on M1M3, and a
  stone under the telescope near the camera.  Inspection with a
  cherry-picker and an infrared camera traced the water to condensation
  on uninsulated Dynalene supply pipes at the Top End Assembly, where
  the dew point exceeded the heat-exchanger return temperature (\rubintkt{SUMMIT-11795},
  \rubintkt{OBS-1479}).

\item[Dec 16--17] Two Dynalene chiller outages, the first coinciding
  with a summit-wide OS patching reboot, caused REBs and CCDs to be
  powered off as an emergency response --- which worked as designed ---
  and put the camera into fault; a few REBs needed a focal-plane
  restart to clear data corruption (\rubintkt{FRACAS-328}).

\item[Dec 22] A changeover between Dynalene chillers warmed the supply
  to the camera and drove REB power-supply board temperatures towards
  their safety limit, forcing HV bias and then the CCDs off for about
  an hour until cooling was restored by switching back (\rubintkt{OBS-1490}).

\item[Dec 21--25] A natural-frequency firmware test on the auxiliary
  telescope's wavefront electronics failed, producing anomalous bias
  levels at both clock rates that no reset sequence would fix.  The
  cause was an inverted-polarity reset signal continuously holding the
  ASPIC in reset; the fix was applied and retested successfully on
  Dec~24, although a new stripe-pattern noise appeared and required
  further study (\rubintkt{SUMMIT-11822}).

\item[Dec 27--29] The science REBs were switched again to the 6.4\,ns
  firmware for an extended holiday calibration campaign, completing a
  partial photon-transfer curve on the night of Dec~27--28 before
  reverting to production firmware; opcode errors recurred on three
  REBs after RCE resets even with idle clear confirmed off, pointing at
  a firmware or boot-sequence cause (\rubintkt{LSSTCAM-153}).

\item[Dec 12, Dec 18] Operations structure changed at year end: the
  camera weekly planning meeting moved from Monday to Thursday, and
  from Dec~28 the summit support scientists moved to a single daily
  shift, with active support from 14:00 to 22:00 local time and
  on-call support thereafter until sunrise.
\end{description}

\paragraph{Sources and method.}
The entries above were extracted from the complete 2025 message
history of the Rubin Observatory Slack channels \#cam-summit-ops and
\#cam-operator, including all thread replies (roughly 2400 and 11\,700
top-level messages respectively, with about 5400 thread replies in
\#cam-summit-ops).  Where the logs give a ticket identifier it is
reproduced verbatim so that the underlying record can be retrieved;
where the logs are ambiguous about a date the entry gives a range.
Statements of cause reflect the assessment recorded in the channel at
the time and are not independent confirmations.
